\introduction
El panorama de la Educación Superior en el siglo XXI está inmerso en una 
dinámica de cambio sin precedentes por rápidos avances tecnológicos y cambios sociales.
La aceleración tecnológica y las fluctuantes demandas del mercado laboral exigen que las instituciones
académicas migren de currículos rígidos hacia planes de estudio dinámicos y
flexibles \citep{unesco2022}. Esta necesidad de flexibilidad introduce un desafío
crítico: la gestión de la variabilidad curricular. La innovación educativa y la gestión 
curricular se han convertido en elementos clave para transformar el paradigma educativo y 
preparar a los estudiantes para un futuro cada vez más digitalizado y globalizado. 

La innovación educativa y la gestión curricular son aspectos fundamentales para el desarrollo de sistemas
educativos actuales, sin embargo, en el presente, resalta la necesidad de una gestión curricular
que promueva la colaboración entre docentes, la adaptación a las necesidades del alumnado
y la actualización de los conteneidos, métodos pedagógicos, planes de estudios, planificación de asignaturas 
e incluso, carreras en general.

La innovación educativa se refiere a la introducción de nuevas ideas, enfoques,
metodologías y tecnologías en el ámbito educativo; busca mejorar la calidad de la enseñanza
y el aprendizaje, fomentando la creatividad, el pensamiento crítico y la resolución de
problemas en los estudiantes. La innovación educativa también promueve la adaptación de
los procesos educativos a las necesidades cambiantes de la sociedad y el entorno laboral.
Por otro lado, la gestión curricular implica el diseño, desarrollo e implementación de planes
de estudio y programas educativos. Por otra parte, la gestión curricular se ocupa de
organizar y estructurar los contenidos, las habilidades y las competencias que se enseñan
en las instituciones educativas; además, implica la evaluación y la mejora continua de los
programas curriculares para asegurar su relevancia y efectividad \citep{gomez2023}.

La innovación educativa y la gestión curricular están estrechamente interrelacionadas.
La introducción de enfoques innovadores en la enseñanza y el aprendizaje requiere una
gestión curricular flexible y adaptable, esto implica revisar y actualizar constantemente los
planes de estudio para integrar nuevas metodologías, tecnologías y recursos educativos.
Igualmente, la gestión curricular efectiva es fundamental para garantizar que los cambios e
innovaciones en el ámbito educativo se implementen de manera coherente y sistemática. La
integración de las \ac{TIC} en el aula permite
ampliar el acceso a la información, fomentar la colaboración entre estudiantes, personalizar
el aprendizaje y desarrollar habilidades digitales clave, de la misma manera, la gestión
curricular puede incluir la planificación y el uso estratégico de la tecnología en los
programas educativos para maximizar su impacto y beneficios \citep{gomez2023}.

En este sentido, la innovación educativa y la gestión curricular deben fundamentarse en una comprensión 
profunda de las particularidades y demandas de los estudiantes. La educación contemporánea coloca en el centro 
la atención a la diversidad, lo que convierte la inclusión y la personalización del aprendizaje en elementos 
indispensables dentro de cualquier propuesta formativa. Esto supone diseñar y organizar planes de estudio flexibles, 
capaces de responder a distintas habilidades, ritmos, estilos cognitivos y necesidades específicas que presentan 
los estudiantes a lo largo de su proceso educativo.

Asimismo, es esencial reconocer que tanto la innovación pedagógica como la gestión curricular no se desarrollan como 
acciones aisladas o independientes, sino que forman parte de un proceso dinámico y articulado. 
Su efectividad depende, en gran medida, del trabajo colaborativo y del compromiso conjunto de todos los actores que 
intervienen en el ámbito educativo. Directivos, docentes, estudiantes, familias y miembros de la comunidad deben 
involucrarse activamente para construir entornos de aprendizaje que favorezcan la equidad, la participación y 
la mejora continua.

De esta manera, la innovación y la gestión curricular no solo buscan transformar prácticas tradicionales, sino también 
crear condiciones que permitan responder de manera pertinente a los desafíos pedagógicos actuales. 
Su verdadero impacto se evidencia cuando logran integrar la voz de los estudiantes, promover la corresponsabilidad 
en la toma de decisiones y generar propuestas educativas que contribuyan al desarrollo integral de cada persona.

La búsqueda de esta flexibilidad ha llevado a la exploración de paradigmas de ingeniería de software aplicados
al ámbito educativo. El concepto de \ac{SPL}
ofrece un enfoque estructurado y probado para gestionar familias de productos
a partir de un conjunto común de activos \citep{apel2013feature}. Cuando este
paradigma se aplica al dominio educativo, surge el concepto de "Líneas de
Productos Educativos", donde un tronco curricular común puede derivar en
múltiples planes de estudio especializados mediante la gestión explícita de su
variabilidad. La investigación en este campo es incipiente pero prometedora;
por ejemplo, estudios en el ecosistema de software educativo \textit{SOLAR} han
demostrado que la modelación de la variabilidad, incluso en componentes
dinámicos como los foros de discusión, es fundamental para crear sistemas
adaptativos que respondan a contextos educativos diversos \citep{guzman2022variabilidad}.

Actualmente, la definición y el mantenimiento de la variabilidad curricular se gestionan 
principalmente mediante documentos estáticos, hojas de cálculo y procesos manuales, un enfoque que resulta 
costoso, lento y altamente propenso a errores. Esta dependencia de herramientas fragmentadas y no integradas 
dificulta la actualización oportuna de la información, genera inconsistencias entre versiones y 
complica la trazabilidad de los cambios realizados. Como consecuencia, pueden diseñarse rutas formativas 
incoherentes o inválidas, afectando la progresión del estudiante y la calidad del proceso educativo. Además, 
este tipo de gestión incrementa la carga de trabajo administrativo y docente, limita la visibilidad global 
del currículo, provoca la divulgación de información desactualizada y compromete la coherencia institucional. 
Todo ello tiene un impacto directo en la satisfacción del estudiante, en los procesos de acreditación y en la 
capacidad de la institución para adaptarse rápidamente a nuevas demandas educativas o del mercado laboral. 
En conjunto, estos efectos negativos evidencian que los métodos tradicionales basados en documentación 
manual ya no son sostenibles, pues ponen en riesgo la calidad, pertinencia y eficiencia del modelo formativo.

La proliferación de modalidades educativas (presenciales, híbridas, en línea,
\ac{MOOC}, \textit{micro-learning}) y enfoques pedagógicos (aprendizaje basado en
proyectos, \textit{flipped classroom}, gamificación) ha creado un panorama de extrema
variabilidad instruccional \citep{bates2019teaching}. Actualmente, esta complejidad se
gestiona predominantemente mediante soluciones \textit{ad-hoc} que resultan en
implementaciones rígidas y de alto costo de mantenimiento \citep{sanchez2020metodologia}.
La educación se enfrenta a la difícil actividad de crear experiencias de
aprendizaje verdaderamente personalizadas para cada estudiante o contexto
institucional sin un esfuerzo manual enorme. Además, cada nuevo curso o
recurso educativo se desarrolla casi desde cero, sin reutilizar componentes
comunes de manera sistemática.

De este análisis, se desprende la siguiente interrogante que define el
\textbf{problema de investigación:}
¿Cómo implementar un servicio \textit{backend} que proporcione la funcionalidad completa
para la gestión, validación y configuración de \ac{FM} en el contexto de Líneas de Productos Educativos?
Tomando como \textbf{objeto de estudio:} 
\ac{SPLE} aplicada al dominio educativo.
Enmarcado en el \textbf{campo de acción:}
servicios \textit{backend} especializados para la gestión de la variabilidad en Líneas de Productos Educativos.
Con el fin de solucionar la problemática planteada, se define como \textbf{objetivo general:}
desarrollar un servicio \textit{backend} para la gestión de la variabilidad
en Líneas de Productos Educativos.

Para dar cumplimiento al objetivo planteado se proponen las siguientes \textbf{tareas investigativas:}
\begin{enumerate}
    \item \textbf{Estudio} exhaustivo del estado del arte en Ingeniería de
    Líneas de Productos de Software, \ac{FM} y su
    aplicación al dominio educativo.
    \item \textbf{Análisis} de diferentes soluciones homólogas para la identificación de características
    generales y tendencias en este tipo de aplicaciones.
    \item \textbf{Diseño} de la arquitectura del servicio \textit{backend}, definiendo los
    componentes principales  y el modelo de datos para
    representar \ac{FM} en el contexto educativo.
    \item \textbf{Selección} de las herramientas, tecnologías y metodología a emplear, para garantizar
    una excelente y eficiente experiencia de desarrollo y alineado con los objetivos del proyecto.
    \item \textbf{Identificación} de los requisitos funcionales y no funcionales necesarios para un
    sistema capaz de gestionar la variabilidad de Líneas de Productos Educativos.
    \item \textbf{Implementación} del servicio \textit{backend} para materializar el sistema que dará respuesta
    a la problemática planteada, asegurando la creación, análisis, validación y configuración de \ac{FM}.
    \item \textbf{Validación} del servicio desarrollado mediante casos de uso representativos
    en el ámbito educativo, evaluando su funcionalidad, rendimiento, utilidad y 
    efectividad para capturar la variabilidad en la educación.
    \item \textbf{Documentación} de los resultados obtenidos y elaboración de recomendaciones para
    futuras investigaciones en el área.
\end{enumerate}

Para darle solución a los objetivos trazados en las tareas de investigación se emplearon los
siguientes \textbf{métodos científicos:}

\begin{itemize}
    \item \textbf{Métodos teóricos:} 
        \begin{itemize}
            \item \textbf{Modelación:} se utilizó para la elaboración de los diagramas del lenguaje de
            modelado y en la representación de las características de la solución.
            \item \textbf{Histórico-Lógico:} Se empleó para comprender la evolución de las técnicas de
            gestión de variabilidad en \ac{SPL}, desde sus orígenes en ingeniería de software
            hasta su aplicación potencial en dominios educativos, analizando su desarrollo
            histórico y lógica interna.
            \item \textbf{Analítico-sintético:} Utilizado en el examen crítico de la bibliografía
            especializada sobre \ac{SPL}, \ac{FM} y sistemas educativos, permitiendo
            sintetizar un marco conceptual adaptado al dominio educativo.
        \end{itemize}
    \item \textbf{Métodos empíricos:}
        \begin{itemize}
            \item \textbf{Observación:} Mediante este método se analizaron herramientas
            existentes de gestión \ac{SPL} (como \textit{pure::variants}, \textit{FeatureIDE}) para
            identificar limitaciones técnicas y requisitos necesarios para el contexto
            educativo específico.
            \item \textbf{Entrevista:} Se entrevista especialistas responsables de diseños
            curriculares que compran aspectos técnicos para identificar los desafíos
            existentes en la flexibilidad y variabilidad de los currículos.
            \item \textbf{Experimentación:} Mediante la implementación de prototipos
            funcionales para validar diferentes enfoques arquitectónicos y algoritmos
            de validación.
        \end{itemize}
\end{itemize}

% \noindent \textbf{Estructura de la investigación:}
% 
% \noindent\textbf{Capítulo 1:} Fundamentación teórica, capítulo encargado de definir los elementos teóricos
% necesarios para el desarrollo de la investigación y los principales conceptos que se emplearán
% durante todo el trabajo. Se realiza un análisis de las soluciones similares y se selecciona la
% metodología de desarrollo de software y las herramientas y tecnologías a utilizar.
% 
% \noindent\textbf{Capítulo 2:} Características y diseño del sistema, capítulo poseedor de los diferentes requisitos
% funcionales y no funcionales de la aplicación a desarrollar, así como de la modelación de la
% solución mediante los artefactos planteados por la metodología de desarrollo a utilizar.
% 
% \noindent\textbf{Capítulo 3:} Implementación y pruebas, capítulo que contiene los estándares de codificación
% utilizados en la implementación, así como las diferentes pruebas a las que será sometida la
% aplicación una vez terminada para verificar su correcto funcionamiento.
